% !TeX root = ../paper.tex

\section{结论}

本文提出 FlexPPI，以主链局部框架和不变点注意力构建对整体旋转和平移保持不变的
几何表示，并从静态多链结构预测逐残基 RMSF、侧链构象熵、SASA 及其不确定性；模型在训练完成后使用 ITrace TCR--pMHC 轨迹进行
外部评估。排除精确 PDB 重叠后，模型在 244 条抗原肽内的平均残基 Spearman 为
0.302，能够恢复部分九肽位置先验和队列平均接触信号，但最活跃残基命中率仅为
34.0\%。完整复合物在全部链上的平均 Spearman 为 0.740，高于逐链 BackFlip-seq
的 0.704、FastProtFlex 的 0.685 和 RMSF-net-PDB 的 0.526，也高于
ProDy-ANM/GNM；但 BackFlip-seq 在线性相关和标准化误差上更好，结构版 BackFlip
在 CDR3 内达到 0.920，超过 FlexPPI 的 0.812。分层结果表明，FlexPPI 的总体排序
优势主要来自非界面、TCR framework 和抗原肽，而不是所有界面区域。

固定坐标伙伴消融揭示了模型利用伙伴信息的方式：相对于只删除对侧伙伴且保留内部
组成的 pMHC/TCR 子复合物，完整输入将 RMSF Spearman 从 0.691 提高至 0.740，
绝对 MAE 从 $0.373\,\text{\AA}$ 降至 $0.346\,\text{\AA}$；全部保留原子的坐标
差为零。加入伙伴同时在接触区域产生强烈且可定位的负向修正，直接接触残基相对
非接触残基额外降低 $0.242\,\text{\AA}$ RMSF、0.225 的 $H_\chi$ 和
$22.87\,\text{\AA}^2$ SASA。该修正在 MHC 接触面方向正确，却错误压低了 TCR 接触面以及
CDR1/2$\beta$ 的柔性，并缩小了 CDR3$\alpha$ 与 CDR3$\beta$ 的链间差异。
因此，FlexPPI 已学到伙伴接触约束，但这一约束过度单向，尚不能稳定刻画界面两侧
及不同 TCR 链的动力学分工。

本文结果同时限定了模型输出的物理解释。共享结合态坐标的删链差值只能表示上下文
响应；缺少独立 apo/holo 轨迹时，不能称为结合诱导动态变化。RMSF、侧链构象熵和
SASA 也不能直接等同于结合自由能。FlexPPI 的当前价值在于以低成本提出残基级动态
假设并暴露静态几何模型的能力边界。完成去同源外部测试、独立配对轨迹验证和构象
采样型基线比较后，方可进一步判断这些动态表示能否为结合响应或亲和力预测提供
可靠增量。
